%%%%%%%%%%%%%%%%%%%%%%%%%%%%%%%%%%%%%%%%%%%%%%%%%%%%%%%%%%%%%%%%%%%%%%%%%%%%%%
%  Unit V appendices
%%%%%%%%%%%%%%%%%%%%%%%%%%%%%%%%%%%%%%%%%%%%%%%%%%%%%%%%%%%%%%%%%%%%%%%%%%%%%%

\startappendices

\chapter{Syllabus-to-Chapter Concordance}

This appendix exists so that coverage can be \emph{audited} rather than assumed.
The left-hand column reproduces the Unit~V syllabus descriptor phrase by phrase,
in the order the official document prints it. The right-hand column gives the
chapter and section that carries it. Every phrase resolves to a location, and no
numbered chapter of this volume contains material that is not traceable to a
phrase in this table.

\begin{mmlongtable}{Concordance for Unit V}{tab:concordance5}%
{@{}P{0.36\textwidth}P{0.57\textwidth}@{}}
\rowcolor{ocre}\thd{Syllabus phrase} & \thd{Where it is covered}\\
\midrule
\endfirsthead
\rowcolor{ocre}\thd{Syllabus phrase} & \thd{Where it is covered}\\
\midrule
\endhead
\multicolumn{2}{@{}l}{\bfseries\color{ocredark} First half --- Conversion Optimization}\\
\midrule
\rlab{What is AIDAS and its role} & Chapter~1 entire. What AIDAS stands for
\S\,1.1; the five stages \S\,1.2; why Satisfaction is critical \S\,1.3; the role
of AIDAS in CRO \S\,1.4; structuring a landing page with AIDAS \S\,1.5.\\
\rowcolor{tablealt}
\rlab{website optimization} & Chapter~2 entire. Conversion, conversion rate and
CRO defined \S\,2.1; why CRO matters \S\,2.2; the conversion funnel and its
drop-offs \S\,2.3; KPIs for website optimization \S\,2.4; the optimisation
workflow \S\,2.5.\\
\rlab{what visitors want to see on the website} & Chapter~3 entire. The ten
visitor questions \S\,3.1; the psychology beneath them \S\,3.2.\\
\rowcolor{tablealt}
\rlab{how to optimize key element and increase the effect of landing on a
particular page} & Chapter~4 entire. Headline \S\,4.1; call to action \S\,4.2;
visuals \S\,4.3; forms \S\,4.4; trust signals and social proof \S\,4.5; page
speed and mobile \S\,4.6.\\
\midrule
\multicolumn{2}{@{}l}{\bfseries\color{ocredark} Second half --- Digital Analytics}\\
\midrule
\rlab{Evolution of Digital Analytics} & Chapter~5 entire. Digital analytics
defined and distinguished from web analytics \S\,5.1; influence on business
decisions \S\,5.2; the six eras \S\,5.3.\\
\rowcolor{tablealt}
\rlab{information about end-to-end customer experience} & Chapter~6 entire. What
it means \S\,6.1; why the siloed view fails \S\,6.2; the seven requirements
\S\,6.3; business value \S\,6.4; refining segmentation \S\,6.5.\\
\rlab{analyst's influence on business} & Chapter~7, \S\,7.1 and \S\,7.2 --- the
interpretation chain and the seven modes of influence. The dashboard of \S\,7.5
and the worked scenarios of \S\,7.6 are that influence in practice.\\
\rowcolor{tablealt}
\rlab{role as a change agent} & Chapter~7, \S\,7.3 --- the shift from reporting
function to change agent, what it requires, and the six barriers.\\
\end{mmlongtable}

\begin{keybox}[Two Notes on Coverage]
\textbf{Deliberate overlap with Volume~4.} ``What visitors want to see'' is named
by both the Unit~IV descriptor (under user-centred design) and the Unit~V
descriptor. Volume~4 treats it as a design brief; Chapter~3 here treats it as a
conversion diagnosis and links it explicitly to the AIDAS\ix{AIDAS} stages.

\smallskip
\textbf{No off-syllabus appendix.} As in Volumes~3 and~4, every topic in the
source material for Unit~V maps to a descriptor phrase.
\end{keybox}

\chapter{Previous-Year Question Bank --- Unit V}

Every question in this appendix is reproduced from an actual RVCE Marketing
Management Continuous Internal Evaluation or Semester End Examination paper for
this course. Answers and pointers follow the official schemes of valuation where
those were released.

\section*{Part A --- Two-Mark Questions}
\addcontentsline{toc}{section}{Part A --- Two-Mark Questions}

\begin{enumerate}[leftmargin=2.2em,itemsep=1.1ex]

\item \textbf{What does the AIDAS model stand for in conversion optimization?}
      \pyqr{Improvement CIE, 4 Jun 2025 --- 2 M}\\
      Attention, Interest, Desire, Action, Satisfaction --- the sequence of
      psychological states a prospect passes through from first exposure to
      post-purchase satisfaction. \hfill\emph{See} \S\,1.1.

\item \textbf{Why is the ``Satisfaction'' component important in the AIDAS
      model?} \pyqr{SEE, Jun/Jul 2025 --- 2 M}\\
      Because the sale is not the end of the funnel. Satisfaction determines
      repurchase, retention and advocacy\ix{advocacy} --- and selling to an existing customer
      succeeds 60--70\% of the time against 5--20\% for a new prospect, with
      repeat customers spending up to 67\% more. A dissatisfied customer also
      broadcasts the failure publicly, damaging future acquisition.
      \hfill\emph{See} \S\,1.3.

\item \textbf{What is the role of AIDAS framework in the conversion
      optimization?} \pyqr{SEE, Oct/Nov 2023 --- 2 M}\\
      It acts as a diagnostic and structuring tool --- locating which
      psychological stage is failing, assigning each page section a defined job,
      sequencing the page accordingly, and generating testable optimisation
      hypotheses. \hfill\emph{See} \S\,1.4.

\item \textbf{What is digital analytics?} \pyqr{CIE-III, Jun 2026 --- 2 M}\\
      The collection, measurement, analysis and interpretation of digital data to
      improve business performance and customer experience.
      \hfill\emph{See} \S\,5.1.

\item \textbf{How does digital analytics influence business decisions?}
      \pyqr{CIE-III, Jun 2026 --- 2 M}\\
      It provides customer insights and performance data that help businesses
      improve marketing strategies and customer engagement --- determining budget
      allocation, product development, segment targeting, experience priorities
      and forecasts. \hfill\emph{See} \S\,5.2.

\item \textbf{What is meant by the ``end-to-end customer experience'' in digital
      analytics?} \pyqr{SEE, Jun/Jul 2025 --- 2 M}\\
      Measuring and understanding the complete customer journey\ix{customer journey} across every
      touchpoint and stage --- from first impression\ix{impression} through research, purchase,
      delivery and support to repeat purchase and advocacy --- as one connected
      experience rather than as disconnected channel reports.
      \hfill\emph{See} \S\,6.1.

\item \textbf{State two key elements that visitors commonly expect to see on a
      well-optimized landing page.} \pyqr{Improvement CIE, 4 Jun 2025 --- 2 M}\\
      A clear benefit-led headline that matches the advertisement or link that
      brought them there, and a single prominent call-to-action button. Also
      acceptable: trust signals such as testimonials or ratings; transparent
      pricing. \hfill\emph{See} \S\,3.1 and \S\,4.1.

\item \textbf{What is Website Optimization?} \pyqr{CIE-III, Jun 2026 --- 2 M}\\
      The process of improving website performance, speed, usability\ix{usability} and search
      visibility in order to increase traffic and conversions.
      \hfill\emph{See} \S\,2.1, and Volume~4 \S\,1.4.

\item \textbf{What is the impact of user behavior analysis on website
      optimization?} \pyqr{SEE, Oct/Nov 2023 --- 2 M}\\
      It reveals what users actually do rather than what is assumed ---
      identifying friction, ignored content, invisible calls to action, rage
      clicks and true scroll depth\ix{scroll depth} --- and produces the hypotheses that A/B testing
      then validates. \hfill\emph{See} \S\,2.3, and Volume~4 \S\,5.4.

\end{enumerate}

\section*{Part B --- Eight- and Ten-Mark Questions}
\addcontentsline{toc}{section}{Part B --- Eight- and Ten-Mark Questions}

\begin{enumerate}[leftmargin=2.2em,itemsep=1.4ex]

\item \textbf{Explain the AIDAS model and its role in website conversion
      optimization. How can key elements of a website be optimized to increase
      landing page effectiveness?} \pyqr{Improvement CIE, 4 Jun 2025 --- 10 M}\\
      \emph{Part one} --- define AIDAS and explain all five stages using the
      stage table, then give the five points on its role in CRO.
      \emph{Part two} --- take the key elements in turn (headline, call to action\ix{call to action},
      visuals, forms, trust signals, speed and mobile) and give the specific
      optimisation for each. Close with the AIDAS-structured landing page layout.
      \hfill\emph{See} Chapter~1 and Chapter~4.

\item \textbf{Explain the AIDAS model in detail. Discuss how each stage plays a
      critical role in conversion optimization, using examples to illustrate how
      businesses can enhance customer journeys through this model.}
      \pyqr{Improvement CIE, 28 Aug 2024 --- 10 M}\\
      As above, but weight the answer towards examples per stage.
      \emph{Attention} --- a streaming service's landing page whose entire first
      screen is one headline, one CTA and nothing else. \emph{Interest} --- a
      product page opening with three benefit bullets rather than technical
      specifications. \emph{Desire} --- a marketplace placing star ratings and
      review counts directly above the buy button. \emph{Action} --- a one-tap
      checkout with saved payment details and guest checkout enabled.
      \emph{Satisfaction} --- an order-confirmation and proactive
      delivery-tracking sequence followed by a review request, converting the
      buyer into an advocate. \hfill\emph{See} \S\,1.2 and \S\,1.5.

\item \textbf{Analyze how a company can use the AIDAS model to structure content
      on its product landing page.} \pyqr{SEE, Jun/Jul 2025, Q10b --- 8 M}\\
      The section-by-section table mapping each part of the page to an AIDAS
      stage, closing with the structuring principle that a section with no
      assignable stage should be deleted. \hfill\emph{See} \S\,1.5.

\item \textbf{Discuss what website visitors typically look for when they land on a
      business webpage. Suggest effective design and content strategies to meet
      those expectations.} \pyqr{Improvement CIE, 4 Jun 2025 --- 10 M}\\
      The ten visitor questions with the paired delivery for each, grounded in
      Krug's first law and the scanning behaviour that follows from it.
      \hfill\emph{See} Chapter~3.

\item \textbf{A travel agency has created a landing page for a limited-time offer
      on holiday packages. Despite a high number of clicks on the ads leading to
      the landing page, the conversion rate is disappointingly low. Evaluate the
      possible reasons for the low conversion rate on the landing page. Suggest
      optimizations for key elements such as headlines, call-to-action buttons,
      and visuals.} \pyqr{Improvement CIE, 28 Aug 2024 --- 10 M}\\
      \emph{Diagnose with AIDAS.} \textbf{Attention held} --- the clicks prove the
      advertisement worked, so the failure is after arrival. \textbf{Interest
      failing} --- poor message match\ix{message match} between the advertisement promise and the
      page headline; the offer is below the fold. \textbf{Desire failing} --- no
      testimonials, ratings, certifications or cancellation policy; pricing is
      opaque; no urgency despite the offer being time-limited; generic stock
      imagery. \textbf{Action failing} --- competing calls to action or none
      dominant; a long enquiry form; site navigation letting visitors wander away;
      slow mobile load.

      \emph{Optimisations --- headline:} restate the advertisement promise with
      specifics, for example ``Kerala backwaters, 4 nights from \rupee 18,999 ---
      book by Sunday'', with a sub-headline listing inclusions.
      \emph{CTA:} one dominant high-contrast button repeated above and below the
      fold, active copy (``Reserve my package''), reassurance microcopy (``Free
      cancellation up to 7 days'').
      \emph{Visuals:} real destination and customer photographs instead of stock, a
      short itinerary video, and a countdown timer.
      \emph{Additional:} cut the form to three fields, place testimonials adjacent
      to the CTA, remove navigation from the landing page, and compress images for
      a sub-three-second mobile load.
      \emph{Validate:} A/B test the headline and CTA, and confirm the fix through
      conversion rate, scroll depth and session recordings.
      \hfill\emph{See} \S\,1.4 and Chapter~4.

\item \textbf{Trace the evolution of digital analytics and its growing importance
      in delivering a seamless end-to-end customer experience.}
      \pyqr{Improvement CIE, 4 Jun 2025 --- 10 M}\\
      Give the six eras --- server logs, page-tag web analytics, visitor and goal
      analytics, multi-channel and social analytics, user-centric event-based
      analytics, and unified predictive privacy-first analytics --- with the tools
      and the limitation of each. Then explain the direction of travel (counting
      files \ra{} predicting behaviour; silos \ra{} unified view), and connect it
      to the end-to-end customer experience material, including the
      decision-constellation point and the seven requirements for genuine
      end-to-end measurement. \hfill\emph{See} \S\,5.3 and Chapter~6.

\item \textbf{A company plans to enter a competitive global market. Explain how
      digital analytics and website optimization can help maintain competitive
      advantage.} \pyqr{CIE-III, Jun 2026 --- 10 M}\\
      The official scheme answer expanded with market-by-market insight,
      localisation of language, currency, payment methods and regional page speed,
      and the argument that speed of learning is the durable advantage.
      \hfill\emph{See} \S\,7.6.

\item \textbf{A food delivery app wants to improve customer retention\ix{customer retention}. Apply
      digital analytics concepts to solve the problem.}
      \pyqr{CIE-III, Jun 2026 --- 10 M}\\
      The official scheme answer expanded into the five-step structure: measure
      \ra{} segment \ra{} diagnose \ra{} act \ra{} re-measure against a control
      group. \hfill\emph{See} \S\,7.6.

\item \textbf{Describe the conversion funnel\ix{conversion funnel} in digital marketing. What are the key
      stages and how can business identify and address drop offs in the funnel?}
      \pyqr{SEE, Oct/Nov 2023, Q9a --- 8 M}\\
      The five-stage funnel table with metrics and drop-off causes, the five
      identification methods, and the stage-by-stage remedies.
      \hfill\emph{See} \S\,2.3.

\item \textbf{Explain the concept of social proof and its impact on conversion
      rates. Provide examples of how business can use social proof effectively in
      marketing.} \pyqr{SEE, Oct/Nov 2023, Q9b --- 8 M}\\
      The definition, the impact argument grounded in the 90\%-against-48\% trust
      asymmetry, and the ten types of social proof with usage guidance --- plus the
      warning that fabricated proof destroys more trust than it borrows.
      \hfill\emph{See} \S\,4.5.

\item \textbf{Define Digital Analytics. Describe the key metrics and KPIs used in
      digital analytics. How can these metrics help business measure the success of
      their online campaigns?} \pyqr{SEE, Oct/Nov 2023, Q10a --- 10 M}\\
      Define digital analytics and distinguish it from web analytics. Present the
      KPIs in structured families --- \emph{acquisition} (sessions, users, traffic
      by channel, cost per click, cost per acquisition), \emph{engagement} (bounce\ix{bounce}
      rate, session duration, pages per session, scroll depth), \emph{conversion}
      (conversion rate, cart abandonment, form completion, revenue, average order
      value), \emph{retention} (repeat purchase rate, churn\ix{churn}, cohort retention, net
      promoter score, lifetime value), and \emph{efficiency} (return on advertising
      spend, lifetime value against acquisition cost).

      Then explain how they measure campaign success: they attribute results to
      specific campaigns and channels, allow comparison against targets and prior
      periods, expose which stage of the funnel a campaign fails at, and convert
      marketing activity into financial terms the business can act on.
      \hfill\emph{See} \S\,5.1 and \S\,2.4.

\item \textbf{What are the key performance indicators (KPIs) used to evaluate
      website optimization, and how do they contribute to increased conversions?}
      \pyqr{SEE, Jun/Jul 2025, Q9a --- 8 M}\\
      The twelve KPIs with the specific contribution each makes to conversion.
      \hfill\emph{See} \S\,2.4.

\item \textbf{Analyze how market segmentation can be refined using digital
      analytics data.} \pyqr{SEE, Jun/Jul 2025, Q9b --- 8 M}\\
      The seven refinements, from behavioural and value-based segmentation through
      to predictive segmentation and continuous automatic refresh.
      \hfill\emph{See} \S\,6.5.

\item \textbf{Design a basic digital dashboard that reflects holistic marketing\ix{holistic marketing}
      performance.} \pyqr{SEE, Jun/Jul 2025, Q10a --- 8 M}\\
      The seven-section dashboard with metrics and the decision each section
      supports, plus the design principles --- one screen, trend not just value,
      target alongside actual, segmented, colour for status only, and delete any
      metric that drives no decision. \hfill\emph{See} \S\,7.5.

\item \textbf{Demonstrate how tools like Google Analytics\ix{Google Analytics} and Hotjar can be used to
      develop insights into user behavior and improve marketing decisions.}
      \pyqr{SEE, Jun/Jul 2025, Q7a --- 8 M}\\
      Analytics supplies the \emph{what} --- channel performance, landing-page
      bounce rates, funnel drop-off, converting segments, cost per acquisition.
      Behavioural tools supply the \emph{why} --- heatmaps showing what attracts
      attention, scroll maps showing whether users reach the CTA, click maps
      exposing rage clicks on non-clickable elements, session recordings revealing
      hesitation, and on-page surveys capturing the abandonment reason in the
      user's own words.

      \emph{Combined workflow:} analytics shows the pricing page loses 60\% of
      visitors \ra{} recordings show users repeatedly clicking a static comparison
      table \ra{} the table is made interactive \ra{} an A/B test validates the
      change \ra{} analytics confirms improved funnel conversion. This is the
      analyst-as-change-agent loop applied in practice.
      \hfill\emph{See} \S\,7.3, and Volume~3 \S\,3.2.

\end{enumerate}

\section*{How Unit V Is Examined}
\addcontentsline{toc}{section}{How Unit V Is Examined}

\begin{keybox}[The Pattern Across Papers]
In the Semester End Examination, Unit~V is Questions~9 and~10 with internal
choice --- 16 marks in two eight-mark sub-divisions --- plus two or three Part~A
objective questions. In the Improvement CIEs it dominates alongside Unit~IV.

\smallskip
\textbf{AIDAS is effectively guaranteed.} It has appeared in every Unit~V paper
available, as both a two-mark definition and a ten-mark explanation. The other
high-probability topics are \textbf{landing page optimisation} for a
low-converting page, the \textbf{evolution of digital analytics}, the
\textbf{conversion funnel} and drop-off analysis, \textbf{social proof}, and
\textbf{KPIs and dashboards}.

\smallskip
As in Unit~IV, most Part-B questions are scenarios. Answer them with the AIDAS
diagnostic structure, or with the measure \ra{} diagnose \ra{} hypothesise \ra{}
test \ra{} implement \ra{} re-measure loop --- not with general commentary.
\end{keybox}

\chapter{Glossary of Key Terms}

\begin{mmlongtable}{Key terms for rapid revision}{tab:glossary5}%
{@{}P{0.26\textwidth}P{0.67\textwidth}@{}}
\rowcolor{ocre}\thd{Term} & \thd{One-line meaning}\\
\midrule
\endfirsthead
\rowcolor{ocre}\thd{Term} & \thd{One-line meaning}\\
\midrule
\endhead
\rlab{A/B test} & Splitting traffic between two versions to find the statistically
better one\\
\rowcolor{tablealt}
\rlab{Above the fold} & The area visible without scrolling\\
\rlab{AIDA} & The classical four-stage model: Attention, Interest, Desire,
Action\\
\rowcolor{tablealt}
\rlab{AIDAS} & Attention, Interest, Desire, Action, Satisfaction\\
\rlab{Cart abandonment rate} & Percentage who begin checkout but do not complete
it\\
\rowcolor{tablealt}
\rlab{CDP} & Customer data platform --- unifies customer data into one profile\\
\rlab{Change agent} & An analyst who drives decisions and organisational change,
not just reports\\
\rowcolor{tablealt}
\rlab{Churn rate} & Percentage of customers lost in a period\\
\rlab{Cohort analysis} & Tracking a group who share a start period over time\\
\rowcolor{tablealt}
\rlab{Conversion} & A defined valuable action completed by a visitor\\
\rlab{Conversion funnel} & The staged path from awareness to purchase and
retention\\
\rowcolor{tablealt}
\rlab{Conversion rate} & Conversions divided by total visitors, times 100\\
\rlab{CRO} & Systematically increasing the percentage of visitors who convert,
using data and experiments\\
\rowcolor{tablealt}
\rlab{Decision constellation} & The non-linear network of simultaneous
micro-decisions replacing the classic funnel\\
\rlab{Digital analytics} & Collection, measurement, analysis and interpretation of
digital data to improve business performance and customer experience\\
\rowcolor{tablealt}
\rlab{Drop-off} & The proportion of users lost between two funnel steps\\
\rlab{End-to-end customer experience} & The complete journey across all
touchpoints, measured as one connected experience\\
\rowcolor{tablealt}
\rlab{HiPPO} & ``Highest paid person's opinion'' --- the bias analytics exists to
counter\\
\rlab{Landing page} & A single-purpose page built for one conversion\\
\rowcolor{tablealt}
\rlab{Last-click attribution} & Crediting the final touchpoint only --- the
classic distortion\\
\rlab{LTV : CAC} & Lifetime value against customer acquisition cost --- the core
efficiency ratio\\
\rowcolor{tablealt}
\rlab{Message match} & Consistency between the advertisement or link and the
landing page it leads to\\
\rlab{Multi-touch attribution} & Distributing conversion credit across all
contributing touchpoints\\
\rowcolor{tablealt}
\rlab{NPS} & Net promoter score --- likelihood of recommendation\\
\rlab{RFM analysis} & Segmenting by recency, frequency and monetary value\\
\rowcolor{tablealt}
\rlab{Risk reversal} & A guarantee or free-return policy that removes the buyer's
perceived risk\\
\rlab{Social proof} & People assuming others' actions reflect correct behaviour\\
\rowcolor{tablealt}
\rlab{Vanity metric} & A number that always rises and drives no decision\\
\end{mmlongtable}

\chapter{Framework Quick-Reference Sheet}

Everything in this volume that is a numbered list, in one place, for the night
before the examination.

\begin{mmlongtable}{Framework quick reference}{tab:quickref5}%
{@{}P{0.28\textwidth}P{0.65\textwidth}@{}}
\rowcolor{ocre}\thd{Framework} & \thd{The sequence}\\
\midrule
\endfirsthead
\rowcolor{ocre}\thd{Framework} & \thd{The sequence}\\
\midrule
\endhead
\rlab{AIDAS} & Attention \ra{} Interest \ra{} Desire \ra{} Action \ra{}
Satisfaction\\
\rowcolor{tablealt}
\rlab{AIDAS states} & ``caught my eye'' \ra{} ``might be relevant'' \ra{} ``I want
this'' \ra{} ``I am buying'' \ra{} ``I am glad I did''\\
\rlab{AIDAS diagnostic} & Immediate bounce = Attention; bounce after scroll =
Interest; add-to-cart without checkout = Action; refunds and no repeat =
Satisfaction\\
\rowcolor{tablealt}
\rlab{Satisfaction economics} & Existing customer 60--70\% against new prospect
5--20\%; repeat customers spend up to 67\% more\\
\rlab{AIDAS role in CRO (five points)} & Locates the failure, maps content to
state, sequences the page, generates hypotheses, prevents over-optimising one
stage\\
\rowcolor{tablealt}
\rlab{AIDAS landing page} & Hero \ra{} value proposition \ra{} proof \ra{} offer
and urgency \ra{} conversion block \ra{} reassurance footer\\
\rlab{Conversion rate} & Conversions divided by total visitors, times 100\\
\rowcolor{tablealt}
\rlab{Why CRO matters} & Compounds on future traffic, cheaper than traffic,
improves every channel's ROI, lowers CAC, replaces opinion with evidence\\
\rlab{Conversion funnel} & Awareness \ra{} interest \ra{} desire \ra{} action
\ra{} satisfaction\\
\rowcolor{tablealt}
\rlab{Identifying drop-offs} & Build the funnel \ra{} segment it \ra{} diagnose
qualitatively \ra{} ask the user \ra{} benchmark\\
\rlab{Twelve CRO KPIs} & Conversion rate, bounce, duration and pages, load time,
exit rate, cart abandonment, form completion, CTA click-through, scroll depth,
mobile gap, return and repeat rate, CPA and ROAS\\
\rowcolor{tablealt}
\rlab{CRO workflow} & Measure \ra{} diagnose \ra{} hypothesise \ra{} prioritise
\ra{} test \ra{} implement \ra{} re-measure\\
\rlab{Ten visitor questions} & Right place, what's in it for me, trust, cost, next
step, findability, speed, mobile, payment safety, human contact\\
\rowcolor{tablealt}
\rlab{Headline rules} & Message match, specific and quantified, benefit-led,
short and dominant, plus a sub-headline\\
\rlab{CTA rules} & One dominant, high contrast, active first-person copy, above
the fold and repeated, 44 pixel tap target, reassurance microcopy\\
\rowcolor{tablealt}
\rlab{Form rules} & Drop unused fields, guest checkout, smart defaults, inline
validation, progress indicator, correct mobile keyboard\\
\rlab{Ten types of social proof} & Testimonials, ratings and reviews, user counts,
case studies, client and media logos, expert endorsement, influencer proof,
user-generated content, real-time activity, wisdom of the crowd\\
\rowcolor{tablealt}
\rlab{Six eras of digital analytics} & Server logs \ra{} page-tag web analytics
\ra{} visitor and goal analytics \ra{} multi-channel and social \ra{} user-centric
event-based \ra{} unified, predictive, privacy-first\\
\rlab{Direction of travel} & Counting files \ra{} visits \ra{} conversions \ra{}
understanding people \ra{} predicting behaviour; silos \ra{} unified view\\
\rowcolor{tablealt}
\rlab{End-to-end requirements (seven)} & Unified identity, consistent event
taxonomy, journey analysis, multi-touch attribution, voice-of-customer,
post-purchase data, consent governance\\
\rlab{Segmentation refinements (seven)} & Behavioural, value-based (RFM),
journey-stage, intent-based, channel-affinity, predictive, continuously
refreshed\\
\rowcolor{tablealt}
\rlab{Interpretation chain} & Product design \ra{} tool setup \ra{} data output
\ra{} human interpretation \ra{} business action\\
\rlab{Analyst's influence (seven)} & Sets the agenda, arbitrates, reallocates
money, prioritises the roadmap, quantifies risk, prevents self-deception,
translates\\
\rowcolor{tablealt}
\rlab{Change agent requires} & Business literacy, credibility, storytelling,
relationships, bias to experimentation, follow-through\\
\rlab{Six barriers} & Data distrust, HiPPO, vanity metrics, silos, privacy
constraints, analysis paralysis\\
\rowcolor{tablealt}
\rlab{Dashboard sections (seven)} & Executive summary, acquisition, engagement,
conversion funnel, retention and loyalty, channel deep-dives, technical health\\
\rlab{Dashboard principles} & One screen, show trend, include target, segment,
colour for status only, delete metrics that drive no decision\\
\end{mmlongtable}
